% Kevin
%=============================================
% Puntos 9, 10: Importancia de la optimización
% y limitaciones del modelo.
%=============================================

\begin{frame}{Optimización y Limitaciones del Modelo}

    \begin{columns}[T]

        %===========================================
        % COLUMNA IZQUIERDA: Optimización
        %===========================================
        \column{0.48\textwidth}

        \begin{block}{¿Por qué fue importante optimizar?}
            \scriptsize
            \begin{itemize}
                \item La optimización permitió identificar \textbf{automáticamente} la región del espacio donde el error era mínimo.
                \vspace{0.15cm}
                \item Sin ella, habría que evaluar manualmente infinitas combinaciones de $(x, y, z)$.
                \vspace{0.15cm}
                \item Es el puente entre el \textbf{modelo matemático} y la \textbf{solución física}.
            \end{itemize}
        \end{block}

        \vspace{0.2cm}

        \begin{exampleblock}{Idea Clave}
            \centering
            \scriptsize
            Buscar el mínimo de $E$ equivale a \textbf{encontrar la fuente sísmica}.
        \end{exampleblock}

        %===========================================
        % COLUMNA DERECHA: Limitaciones
        %===========================================
        \column{0.48\textwidth}

        \begin{block}{Limitaciones del Modelo}
            \scriptsize
            \begin{itemize}
                \item \textbf{Medio homogéneo:} se asumió propagación uniforme, sin variaciones del terreno.
                \vspace{0.1cm}
                \item \textbf{Mínimos locales:} la optimización podría converger a un mínimo no global.
                \vspace{0.1cm}
                \item \textbf{Sensibilidad al ruido:} el ruido en las mediciones afecta la precisión de la estimación.
                \vspace{0.1cm}
                \item \textbf{Sensores limitados:} solo 12 estaciones restringen la resolución espacial.
            \end{itemize}
        \end{block}

    \end{columns}

    \vspace{0.15cm}
    \begin{center}
        \footnotesize
        \textbf{Reconocer las limitaciones fortalece la validez científica del modelo.}
    \end{center}

\end{frame}